% !TEX program = xelatex
% !TEX encoding = UTF-8 Unicode
% =========================================================
% 开源单文件中文技术简历模板
% 编译方式：XeLaTeX
% 说明：本文件不依赖项目内自定义 cls/sty/字体目录；仅使用 TeX 发行版常见宏包。
%       示例内容均为虚构信息，请按个人真实经历替换。
% =========================================================

\documentclass[10pt,a4paper]{article}

% ---------- Page / packages ----------
\usepackage[a4paper,left=1.35cm,right=1.35cm,top=1.00cm,bottom=1.05cm]{geometry}
\usepackage{fontspec}
\usepackage{xeCJK}
\usepackage{xcolor}
\usepackage{enumitem}
\usepackage{tabularx}
\usepackage{array}
\usepackage{fontawesome5}
\usepackage{needspace}
\usepackage[colorlinks=false,pdfborder={0 0 0}]{hyperref}

% ---------- Fonts: robust fallbacks ----------
\IfFontExistsTF{Times New Roman}{%
  \setmainfont{Times New Roman}
}{%
  \IfFontExistsTF{Tinos}{\setmainfont{Tinos}}{\setmainfont{TeX Gyre Termes}}
}
\IfFontExistsTF{Noto Serif CJK SC}{%
  \setCJKmainfont[BoldFont=Noto Serif CJK SC Bold]{Noto Serif CJK SC}
  \setCJKsansfont[BoldFont=Noto Sans CJK SC Bold]{Noto Sans CJK SC}
  \setCJKfamilyfont{resumehei}[BoldFont=Noto Sans CJK SC Bold]{Noto Sans CJK SC}
  \setCJKfamilyfont{resumekai}[AutoFakeBold=1.2]{FandolKai-Regular.otf}
}{%
  \setCJKmainfont[AutoFakeBold=2.0]{FandolSong-Regular.otf}
  \setCJKsansfont[AutoFakeBold=2.0]{FandolHei-Regular.otf}
  \setCJKfamilyfont{resumehei}[AutoFakeBold=2.0]{FandolHei-Regular.otf}
  \setCJKfamilyfont{resumekai}[AutoFakeBold=1.2]{FandolKai-Regular.otf}
}
\newcommand{\resumehei}{\CJKfamily{resumehei}}
\newcommand{\resumekai}{\CJKfamily{resumekai}}

% ---------- Colors ----------
\definecolor{ResumeBlack}{RGB}{22,22,22}
\definecolor{ResumeGray}{RGB}{92,92,92}
\definecolor{LightGray}{RGB}{238,238,238}

% ---------- Global layout ----------
\pagestyle{empty}
\setlength{\parindent}{0pt}
\setlength{\parskip}{0pt}
\setlength{\tabcolsep}{0pt}
\renewcommand{\arraystretch}{1.02}
\linespread{1.02}
\emergencystretch=4em
\hbadness=10000
\XeTeXlinebreaklocale "zh"
\XeTeXlinebreakskip = 0pt plus 0.8pt

\setlist[itemize]{
  leftmargin=1.45em,
  labelsep=0.42em,
  itemsep=1.2pt,
  topsep=1.2pt,
  parsep=0pt,
  partopsep=0pt
}
\setlist[enumerate]{
  leftmargin=2.05em,
  labelsep=0.42em,
  itemsep=2pt,
  topsep=2pt,
  parsep=0pt,
  partopsep=0pt
}

% ---------- Components ----------
\newcommand{\ResumeName}[1]{%
  \begin{center}
    {\fontsize{22pt}{25pt}\selectfont\bfseries #1}\par
    \vspace{0.52em}
  \end{center}
}

\newcommand{\ContactLine}{%
  \begin{center}
    {\fontsize{9.7pt}{11pt}\selectfont
    \faPhone\enspace (+86)138-0000-0000
    \hspace{2.0em}
    \href{mailto:nienie@example.com}{\faEnvelope\enspace nienie@example.com}
    \hspace{2.0em}
    \href{https://github.com/}{\faGithub\enspace @nienie}}
  \end{center}
  \vspace{0.22em}
}

\newcommand{\ResumeSection}[2]{%
  \Needspace{3.5\baselineskip}%
  \vspace{0.68em}%
  {\fontsize{13pt}{15pt}\selectfont\bfseries\textcolor{ResumeBlack}{#1\hspace{0.42em}{\resumehei #2}}}\par
  \vspace{-0.28em}%
  \noindent\rule{\linewidth}{0.55pt}\par
  \vspace{0.28em}%
}

\newcommand{\DatedLine}[2]{%
  \noindent\begin{tabularx}{\linewidth}{@{}>{\raggedright\arraybackslash}X r@{}}
    #1 & {\textcolor{ResumeBlack}{#2}}
  \end{tabularx}\par\vspace{-0.10em}%
}

\newcommand{\CompanyBlock}[3]{%
  \Needspace{4\baselineskip}%
  \vspace{0.18em}%
  \noindent\begin{tabularx}{\linewidth}{@{}>{\raggedright\arraybackslash}X r@{}}
    {\bfseries #1}\quad{\resumekai\textcolor{ResumeGray}{#2}} & {\bfseries #3}
  \end{tabularx}\par
  \vspace{0.02em}%
}

\newcommand{\CompanySep}{%
  \vspace{0.42em}%
  \noindent\textcolor{LightGray}{\rule{\linewidth}{0.40pt}}\par
  \vspace{0.36em}%
}

\newcommand{\ProjectTitle}[1]{%
  \vspace{0.34em}%
  \noindent{\bfseries #1}\par
  \vspace{-0.14em}%
}

\newenvironment{ResumeBullets}{%
  \begin{itemize}[
    leftmargin=1.45em,
    labelsep=0.42em,
    itemsep=1.2pt,
    topsep=1.1pt,
    parsep=0pt,
    partopsep=0pt
  ]
}{%
  \end{itemize}\vspace{0.02em}
}

\newenvironment{InternBullets}{%
  \fontsize{10.1pt}{12.35pt}\selectfont
  \begin{itemize}[
    leftmargin=1.45em,
    labelsep=0.42em,
    itemsep=2.2pt,
    topsep=1.8pt,
    parsep=0.6pt,
    partopsep=0pt
  ]
}{%
  \end{itemize}\vspace{0.10em}
}

\newcommand{\ResearchSubTitle}[1]{%
  \Needspace{3\baselineskip}%
  \vspace{0.42em}%
  \noindent{\bfseries\resumehei #1}\par
  \vspace{-0.10em}%
}

\newcommand{\CoFirst}{\textsuperscript{\ensuremath{\dagger}}}
\newcommand{\Corr}{\textsuperscript{\ensuremath{*}}}
\newcommand{\PubLevel}[1]{\textbf{#1}}
\newcommand{\PubStatus}[1]{\textit{#1}}

\begin{document}
\fontsize{10.1pt}{11.85pt}\selectfont

\ResumeName{捏捏}
\ContactLine

\ResumeSection{\faGraduationCap}{\resumekai 教育背景}
\DatedLine{\textbf{星河大学}，计算机科学与技术，{\resumekai 硕士研究生}}{2026.09 -- 2029.06}
\DatedLine{\textbf{云岚大学}，软件工程，{\resumekai 工学学士}}{2022.09 -- 2026.06}
\begin{ResumeBullets}
  \item \textbf{学业表现：}专业排名 1/120，GPA 3.92/4.00；获校级一等奖学金、优秀学生干部、优秀毕业生等荣誉
  \item \textbf{核心课程：}数据结构、操作系统、计算机网络、数据库系统、机器学习、分布式系统、编译原理
  \item \textbf{英语能力：}CET-6；具备英文论文阅读、技术文档写作与跨团队技术沟通能力
\end{ResumeBullets}

\ResumeSection{\faBriefcase}{\resumekai 实习经历}

\CompanyBlock{云脉科技 | CloudPulse Labs}{AI 基础设施研发实习}{2025.06 -- 2025.09}
\ProjectTitle{\normalfont 大模型推理服务构建与性能优化}
\begin{InternBullets}
  \item \textbf{背景与问题：}团队内部多模型推理服务在模型加载、请求排队、日志追踪与灰度发布方面缺少统一流程，导致服务上线成本高、问题定位链路长、不同模型版本之间的性能差异难以稳定复现。
  \item \textbf{工作与成果：}参与推理服务框架建设，协助实现模型配置管理、批量请求调度、推理日志采集与基础监控能力；整理多类压测场景并沉淀性能分析报告，帮助团队识别显存占用、首 token 延迟和并发吞吐等关键瓶颈，提升模型服务迭代效率。
\end{InternBullets}

\ProjectTitle{\normalfont 离线评测 Pipeline 与实验复现治理}
\begin{InternBullets}
  \item \textbf{背景与问题：}模型效果评估依赖分散脚本和人工汇总，存在数据版本不统一、评测指标口径不一致、实验结果难以追溯等问题，影响模型选型与上线决策效率。
  \item \textbf{工作与成果：}参与搭建离线评测 pipeline，统一数据集组织、Prompt 模板管理、批量推理与结果聚合流程；协助完成多轮模型版本、参数配置和提示词策略对比实验，沉淀可复用评测模板与失败样例归因记录，降低后续实验复现成本。
\end{InternBullets}

\CompanySep

\CompanyBlock{青竹智能 | BambooAI}{算法工程实习}{2024.10 -- 2025.01}
\ProjectTitle{\normalfont 多场景文本生成质量评估与提示词优化}
\begin{InternBullets}
  \item \textbf{背景与问题：}业务侧文本生成需求覆盖摘要、问答、改写和结构化抽取等场景，不同场景对事实一致性、格式稳定性与响应风格要求差异较大，人工评估成本较高。
  \item \textbf{工作与成果：}协助整理多类业务样例、构建评测规则与人工标注参考，参与 Prompt 模板迭代、坏例归因和生成结果对比分析；通过统一样例组织和评估口径，使模型效果讨论更加可复现，为后续模型选型和产品化落地提供依据。
\end{InternBullets}

% \clearpage

\ResumeSection{\faGithub}{\resumekai 开源经历}
\begin{ResumeBullets}
  \item \textbf{工程规范实践：}熟悉 Git 协作、Issue/PR 流程、单元测试、CI 检查与代码 Review，能够在多人协作项目中保持稳定交付
\end{ResumeBullets}

\ResumeSection{\faTrophy}{\resumekai 竞赛获奖}
\DatedLine{\textbf{全国大学生软件设计竞赛，}国家级一等奖}{2025.12}

\ResumeSection{\faFile}{\resumekai 科研经历}
\vspace{-0.34em}
{\footnotesize\resumekai 注：\CoFirst{} 表示共同第一作者，\Corr{} 表示通讯作者；以下论文条目为模板示例。}\par

\ResearchSubTitle{代表性科研成果}
\begin{enumerate}[label={\normalfont[\arabic*]},series=pubs]
  \item \textbf{Nienie}\CoFirst, Ming Chen\CoFirst, Rui Zhao, and Lan Xu\Corr.
  Adaptive Context Routing for Efficient Multi-Agent Large Language Model Serving.
  Submitted to \textit{International Conference on Machine Learning Systems}, 2026
  （\PubLevel{ CCF-B}, \PubStatus{Under Review}）.
\end{enumerate}

\ResearchSubTitle{合作科研成果}
\begin{enumerate}[label={\normalfont[\arabic*]},resume=pubs]
  \item Yue Ma, Kai Huang, \textbf{Nienie}, and Zhen Liu\Corr.
  Resource-Aware Model Deployment for Edge-Cloud Collaborative Intelligence.
  In \textit{International Conference on Distributed Computing Systems}, 2025（\PubLevel{CCF-B}）.
\end{enumerate}

\end{document}
